% ====================================================================
\chapter{Absolute Value and the Triangle Inequality}\label{ch:abs}
% ====================================================================

\begin{chapterintro}
The absolute value of a number measures its size regardless of sign; the absolute value of a difference
measures the distance between two numbers. Distance is the key idea behind limits: a sequence approaches a
number when the distance between them becomes as small as we please. This chapter defines the absolute
value, proves its basic properties, establishes the triangle inequality, and shows how to solve inequalities
involving absolute values.
\end{chapterintro}

\begin{prerequisites}
The properties of the real numbers and the sign rules (\cref{ch:numbers}, especially \cref{prop:signs}); proof by
cases (\cref{sec:proof-types}); solving linear and product inequalities (\cref{ch:ineq}); intervals
(\cref{sec:intervals}).
\end{prerequisites}

\section{Definition}\label{sec:abs-def}

\begin{definition}[Absolute value][def:abs]
The \term{absolute value} $|x|$ of a real number $x$ is defined as follows:
\[
|x|=\begin{cases} x & \text{if } x\ge 0,\\ -x & \text{if } x<0.\end{cases}
\]
Equivalently, $|x|=\max\{x,-x\}$, the larger of the two numbers $x$ and $-x$. \src{L2, slide 10}
\end{definition}

For a non-negative number the absolute value changes nothing; for a negative number it removes the minus sign.
The two descriptions agree: if $x\ge0$ then $x\ge -x$ (since $-x\le0\le x$), so $\max\{x,-x\}=x$; if $x<0$ then
$-x>0>x$, so $\max\{x,-x\}=-x$.

\begin{example}[Computing absolute values]
$|2|=2$, \ $|-5|=5$, \ $\big|\tfrac73\big|=\tfrac73$, \ $|-\sqrt{11}|=\sqrt{11}$, \ $|0|=0$. \src{L2, slide 10}
\end{example}

\begin{warning}[$|{-x}|$ is not always $x$]
It is tempting to think ``the absolute value removes the minus sign'', so that $|{-x}|=x$. That is true only if
$x\ge0$. For $x=-3$, $|{-x}|=|3|=3=-x$. The correct statement is $|{-x}|=|x|$ (\cref{prop:abs-props}).
\end{warning}

\begin{figure}[ht]
\centering
\begin{tikzpicture}
\begin{axis}[stat, width=8cm, height=5.2cm, xmin=-3.3, xmax=3.3, ymin=-0.3, ymax=3.4, axis x line=middle,
  axis y line=middle, xlabel={$x$}, ylabel={$y$}, x label style={anchor=west}, y label style={anchor=south},
  xtick={-3,-2,-1,1,2,3}, ytick={1,2,3}]
  \addplot[curve, accent, domain=-3.1:0] {-x};
  \addplot[curve, accent, domain=0:3.1] {x};
  \node[lab, accent] at (axis cs:2.2,3.0) {$y=|x|$};
\end{axis}
\end{tikzpicture}
\caption{The graph of $y=|x|$. To the right of $0$ it coincides with $y=x$; to the left it coincides with
$y=-x$. The graph has a corner at the origin.}
\label{fig:abs-graph}
\end{figure}

\section{Properties}\label{sec:abs-props}

\begin{proposition}[Properties of the absolute value][prop:abs-props]
For all real numbers $a$ and $b$: \src{L2, slide 10}
\begin{enumerate}[label=(\roman*)]
  \item $|a|\ge a$;
  \item $|a|\ge -a$;
  \item $|{-a}|=|a|$;
  \item $|a\cdot b|=|a|\cdot|b|$;
  \item if $b\ne0$, then $\left|\dfrac ab\right|=\dfrac{|a|}{|b|}$ (this follows from the previous property).
\end{enumerate}
In addition, $|a|\ge0$, with $|a|=0$ if and only if $a=0$.
\end{proposition}
\emph{Proof.} (i) and (ii) are immediate from $|a|=\max\{a,-a\}$: the maximum of two numbers is at least as large
as each of them. Combining them with \cref{prop:signs} gives the useful double inequality
\begin{equation}\label{eq:abs-sandwich}
-|a|\le a\le |a|.
\end{equation}
(iii) $|{-a}|=\max\{-a,-(-a)\}=\max\{-a,a\}=|a|$.

(iv) Consider the signs. If $a\ge0$ and $b\ge0$ then $ab\ge0$ and $|ab|=ab=|a||b|$. If $a\ge0$ and $b<0$ then
$ab\le0$, so $|ab|=-ab=a\cdot(-b)=|a||b|$. The case $a<0$, $b\ge0$ is symmetric. If $a<0$ and $b<0$ then $ab>0$,
so $|ab|=ab=(-a)(-b)=|a||b|$.

(v) Apply (iv) to $a=\frac ab\cdot b$: $|a|=\big|\frac ab\big|\cdot|b|$. Since $b\ne0$, $|b|>0$, and dividing by
$|b|$ gives the result.

Finally, if $a>0$ then $|a|=a>0$; if $a<0$ then $|a|=-a>0$; and $|0|=0$. \hfill$\square$

\begin{external}[Absolute value and square roots]
For every real $a$, $|a|^2=a^2$ and $\sqrt{a^2}=|a|$ (where $\sqrt{\ }$ denotes the non-negative square root).
The second identity is a frequent source of errors: $\sqrt{(-3)^2}=\sqrt9=3$, not $-3$.
\end{external}

\section{Absolute value as distance}\label{sec:abs-distance}

\begin{core}[Distance on the number line]
For real numbers $x$ and $y$, $|x-y|$ is the \term{distance} between $x$ and $y$ on the number line. In
particular $|x|=|x-0|$ is the distance from $x$ to $0$. \src{L2, slide 11}
\end{core}

Distance does not depend on direction: $|x-y|=|y-x|$ by \cref{prop:abs-props}(iii). For example, the distance
between $-2$ and $5$ is $|-2-5|=7=|5-(-2)|$.

Statements about distance translate directly into absolute values, and vice versa:
\begin{center}
\begin{tabular}{@{}ll@{}}
\toprule
In words & In symbols\\
\midrule
$x$ is within $2$ of $3$ & $|x-3|<2$\\
$x$ is at least $4$ away from $-1$ & $|x+1|\ge4$\\
$a_n$ is within $\varepsilon$ of $x$ & $|a_n-x|<\varepsilon$\\
\bottomrule
\end{tabular}
\end{center}
The last line is exactly the condition that appears in the definition of a limit (\cref{ch:limits}).

\section{The triangle inequality}\label{sec:triangle}

\begin{theorem}[Triangle inequality][thm:triangle3]
For any real numbers $x,y,z$,
\[
|x-z|\le|x-y|+|y-z|. \tag*{\srcin{L2, slide 11; HW-1 Problem B1}}
\]
\end{theorem}

\begin{intuition}
$|x-z|$ is the distance from $x$ to $z$; $|x-y|$ is the distance from $x$ to $y$; $|y-z|$ is the distance from
$y$ to $z$. Going from $x$ to $z$ directly should not be longer than going first from $x$ to $y$ and then from
$y$ to $z$. A detour cannot be shorter than the direct route. The name comes from geometry: in a triangle, one
side is never longer than the sum of the other two.
\end{intuition}

\begin{sourcenote}[The proof]
The slides pose the triangle inequality as Problem B1 of the first homework and refer to ``notes'' for the
argument, but no proof appears in the materials provided. The proof below is supplied for completeness.
\end{sourcenote}

The cleanest route is to prove a two-number version first.

\begin{theorem}[Triangle inequality, two-number form][thm:triangle2]
For any real numbers $a$ and $b$, $|a+b|\le|a|+|b|$.
\end{theorem}
\emph{Proof.} By \eqref{eq:abs-sandwich}, $-|a|\le a\le|a|$ and $-|b|\le b\le|b|$. Adding these (\cref{ch:numbers},
Question 3.5 applied to $\le$),
\[
-(|a|+|b|)\le a+b\le |a|+|b|.
\]
Now $|a+b|$ equals either $a+b$ or $-(a+b)$. The right-hand inequality gives $a+b\le|a|+|b|$; the left-hand
inequality, multiplied by $-1$ (reversing it), gives $-(a+b)\le|a|+|b|$. In both cases $|a+b|\le|a|+|b|$.
\hfill$\square$

\emph{Proof of \cref{thm:triangle3}.} Put $a=x-y$ and $b=y-z$. Then $a+b=x-z$, and \cref{thm:triangle2} gives
$|x-z|\le|x-y|+|y-z|$. \hfill$\square$

The slides ask how the three-number form can be used to prove two further inequalities. Both follow by choosing
the three numbers cleverly.

\begin{example}[Applying the triangle inequality][ex:triangle-apply]
Use \cref{thm:triangle3} to show that for any real $a,b$: (i) $|a+b|\le|a|+|b|$; (ii) $|a|\le|a-b|+|b|$.
\src{L2, slide 11}
\solution
(i) Take $x=a$, $y=0$, $z=-b$. Then $|x-z|=|a-(-b)|=|a+b|$, $|x-y|=|a-0|=|a|$ and $|y-z|=|0-(-b)|=|b|$. So $|a+b|\le|a|+|b|$.

(ii) Take $x=a$, $y=b$, $z=0$. Then $|x-z|=|a|$, $|x-y|=|a-b|$ and $|y-z|=|b|$. So $|a|\le|a-b|+|b|$.
\end{example}

\begin{external}[Reverse triangle inequality]
Rearranging (ii) gives $|a|-|b|\le|a-b|$; exchanging $a$ and $b$ gives $|b|-|a|\le|a-b|$. Together:
$\big||a|-|b|\big|\le|a-b|$. The sizes of two numbers differ by no more than the distance between them.
\end{external}

\begin{keyformula}[The triangle inequality][kf:triangle]{|x-z|\le|x-y|+|y-z|\qquad\text{equivalently}\qquad |a+b|\le|a|+|b|}
\fsymbols $x,y,z,a,b$ are any real numbers; $|\cdot|$ is the absolute value.
\fmeaning The distance between two points is at most the length of any route between them that passes through a
third point. In the second form: the size of a sum is at most the sum of the sizes.
\forigin From $-|a|\le a\le|a|$ and $-|b|\le b\le|b|$ by adding (\cref{thm:triangle2}); the three-number form
follows by substituting $a=x-y$, $b=y-z$.
\fuse To bound the distance between two numbers by passing through a convenient intermediate number; to bound
the size of a sum; above all in proofs about limits (\cref{ex:no-limit}).
\fexample With $x=1$, $y=x_0$, $z=0$: $1=|1-0|\le|1-x_0|+|x_0-0|$. If both distances on the right were less than
$\tfrac14$, we would get $1<\tfrac12$, which is false. This is exactly the argument that $1,0,1,0,\dots$ has no
limit.
\fmistakes Writing $|a+b|=|a|+|b|$ (equality holds only when $a$ and $b$ do not have opposite signs: $|3+(-5)|=2$
but $|3|+|-5|=8$); writing $|a-b|\le|a|-|b|$ (false: $a=0$, $b=1$ gives $1\le-1$).
\frearrange $|a|-|b|\le|a-b|$ and $\big||a|-|b|\big|\le|a-b|$ (reverse triangle inequality); $|a-b|\le|a|+|b|$
(apply the two-number form to $a$ and $-b$).
\fchanges Equality $|a+b|=|a|+|b|$ holds exactly when $ab\ge0$. In the three-number form, equality holds exactly
when $y$ lies between $x$ and $z$ (inclusive): moving $y$ outside that interval makes the detour strictly longer.
\fexam As a step in a proof (``by the triangle inequality, \dots''), particularly in arguments by contradiction
about limits; or as a short ``show that'' question.
\end{keyformula}

\section{Solving inequalities with absolute values}\label{sec:abs-ineq}

\begin{proposition}[Within a distance][prop:abs-within]
Let $a\in\R$ and $r>0$. Then
\begin{align*}
|x-a|<r\quad&\Longleftrightarrow\quad a-r<x<a+r\quad\Longleftrightarrow\quad x\in(a-r,\,a+r),\\
|x-a|>r\quad&\Longleftrightarrow\quad x<a-r\ \text{ or }\ x>a+r.
\end{align*}
The same holds with $\le$ and $\ge$, with the endpoints included.
\end{proposition}
\emph{Proof.} Split by the sign of $x-a$. If $x\ge a$, then $|x-a|=x-a$, and $|x-a|<r\Leftrightarrow x<a+r$; combined
with $x\ge a$ this gives $a\le x<a+r$. If $x<a$, then $|x-a|=a-x$, and $|x-a|<r\Leftrightarrow x>a-r$; combined with
$x<a$ this gives $a-r<x<a$. The union of the two cases is $a-r<x<a+r$. The second statement is the negation of
the first with $\le$ (\cref{sec:negation}): $\text{not}(a-r\le x\le a+r)\Leftrightarrow x<a-r$ or $x>a+r$.
\hfill$\square$

\begin{figure}[ht]
\centering
\begin{tikzpicture}[x=10mm]
  \draw[-{Stealth}] (-0.8,0) -- (7.8,0);
  \draw[accent, line width=2.4pt] (1.5,0) -- (5.5,0);
  \node[circle, draw=accent, fill=white, thick, inner sep=1.5pt] at (1.5,0) {};
  \node[circle, draw=accent, fill=white, thick, inner sep=1.5pt] at (5.5,0) {};
  \draw (3.5,0.12) -- (3.5,-0.12);
  \node[below, font=\small] at (1.5,-0.1) {$a-r$}; \node[below, font=\small] at (3.5,-0.1) {$a$};
  \node[below, font=\small] at (5.5,-0.1) {$a+r$};
  \draw[decorate, decoration={brace, amplitude=4pt}] (3.5,0.25) -- (5.5,0.25) node[midway, above=4pt, font=\small] {$r$};
  \draw[decorate, decoration={brace, amplitude=4pt}] (1.5,0.25) -- (3.5,0.25) node[midway, above=4pt, font=\small] {$r$};
\end{tikzpicture}
\caption{The set $\{x\in\R\mid |x-a|<r\}$ is the open interval of radius $r$ centred at $a$: all points whose
distance from $a$ is less than $r$.}
\label{fig:abs-interval}
\end{figure}

\begin{example}[Inequalities with absolute values]
Solve (a) $|x-3|<2$; (b) $|2x+1|\le5$; (c) $|x|>3$; (d) $|x-1|<|x+2|$.
\solution
(a) By \cref{prop:abs-within} with $a=3$, $r=2$: $1<x<5$, so $x\in(1,5)$. In words: the numbers within distance
$2$ of $3$.

(b) $|2x+1|\le5\Leftrightarrow-5\le2x+1\le5\Leftrightarrow-6\le2x\le4\Leftrightarrow-3\le x\le2$, so $x\in[-3,2]$.
(Alternatively write $|2x+1|=2|x+\tfrac12|$ and use distance from $-\tfrac12$.)

(c) $x<-3$ or $x>3$: $x\in(-\infty,-3)\cup(3,\infty)$.

(d) Both sides are non-negative, and for non-negative numbers $u<v\Leftrightarrow u^2<v^2$ (by \cref{ch:numbers}, Question 3.7
and its converse). So the inequality is equivalent to $(x-1)^2<(x+2)^2$, i.e.\ $x^2-2x+1<x^2+4x+4$, i.e.\
$-6x<3$, i.e.\ $x>-\tfrac12$. In words: $x$ is closer to $1$ than to $-2$ exactly when it lies to the right of the
midpoint $-\tfrac12$.
\end{example}

\begin{warning}[Splitting $|x-a|<r$ wrongly]
$|x-3|<2$ does \emph{not} mean ``$x-3<2$'' alone; that drops the lower bound. Nor does $|x|>3$ mean $-3>x>3$, which
is impossible. Inequalities of the form $|\cdot|<r$ give a single interval (``and''); those of the form $|\cdot|>r$
give two pieces (``or'').
\end{warning}

\begin{example}[An economic tolerance band]
A central bank targets an inflation rate of $2\%$ and regards outcomes within one percentage point of the target as
acceptable. Write the acceptable range as an absolute-value inequality and as an interval.
\solution
If $\pi$ denotes inflation in per cent, the condition is $|\pi-2|\le1$, i.e.\ $1\le\pi\le3$, or $\pi\in[1,3]$.
Inflation outside this band, $|\pi-2|>1$, means $\pi<1$ or $\pi>3$.
\end{example}

% --------------------------------------------------------------------
\section{Chapter review}
% --------------------------------------------------------------------
\subsection*{Summary}
$|x|$ equals $x$ for $x\ge0$ and $-x$ for $x<0$; equivalently $|x|=\max\{x,-x\}$. It is never negative and is zero
only at $0$. It satisfies $|a|\ge a$, $|a|\ge-a$, $|{-a}|=|a|$, $|ab|=|a||b|$ and $|a/b|=|a|/|b|$. The quantity
$|x-y|$ is the distance between $x$ and $y$. The triangle inequality, $|x-z|\le|x-y|+|y-z|$, says a detour is never
shorter than the direct route; its two-number form is $|a+b|\le|a|+|b|$. Inequalities $|x-a|<r$ describe intervals of
radius $r$ about $a$, and $|x-a|>r$ describes the two pieces outside.

\subsection*{Key definitions}
\begin{itemize}
  \item $|x|=x$ if $x\ge0$, $|x|=-x$ if $x<0$; equivalently $|x|=\max\{x,-x\}$.
  \item $|x-y|$: the distance between $x$ and $y$.
\end{itemize}

\subsection*{Key equations}
\begin{gather*}
-|a|\le a\le|a|,\qquad |{-a}|=|a|,\qquad |ab|=|a||b|,\qquad \Big|\frac ab\Big|=\frac{|a|}{|b|}\ (b\ne0),\\
|x-z|\le|x-y|+|y-z|,\qquad |a+b|\le|a|+|b|,\qquad |a|\le|a-b|+|b|,\\
|x-a|<r\ \Leftrightarrow\ a-r<x<a+r.
\end{gather*}

\subsection*{Connections}
The proofs use cases (\cref{sec:proof-types}) and the sign rules of \cref{prop:signs}. Absolute-value inequalities
reduce to the linear inequalities of \cref{ch:ineq}. The distance interpretation and the triangle inequality are
the main tools of \cref{ch:limits}, where they are used to prove that a sequence converges or that it does not.

\subsection*{Common mistakes}
\begin{itemize}
  \item Writing $|{-x}|=x$ or $\sqrt{x^2}=x$ without the condition $x\ge0$.
  \item Writing $|a+b|=|a|+|b|$ or $|a-b|=|a|-|b|$.
  \item Solving $|x-a|<r$ as $x-a<r$ only.
  \item Squaring both sides of an inequality whose sides might be negative.
\end{itemize}

\subsection*{Exam checklist}
\begin{checklist}
  \item State the definition of $|x|$ in both forms.
  \item State and prove the properties $|a|\ge a$, $|{-a}|=|a|$, $|ab|=|a||b|$.
  \item State the triangle inequality and explain it as a statement about distances.
  \item Derive $|a+b|\le|a|+|b|$ and $|a|\le|a-b|+|b|$ from the three-number form.
  \item Solve inequalities of the forms $|x-a|<r$, $|x-a|\ge r$ and $|\,\cdot\,|<|\,\cdot\,|$.
\end{checklist}

% --------------------------------------------------------------------
\section{Practice questions}
% --------------------------------------------------------------------
\pq[basic] Evaluate $|-7|$, $|3-8|$, $|2-\sqrt5|$, $|-\tfrac{4}{9}|$ and $|x^2+1|$ (for real $x$).

\pq[basic] Solve (a) $|x+4|<1$; (b) $|3x-6|\ge 9$; (c) $|5-x|\le 0$.

\pq[conceptual] Give an example of real numbers $a,b$ with $|a+b|<|a|+|b|$, and one with $|a+b|=|a|+|b|$. What
distinguishes the two situations?

\pq[mathematical] Prove that for all real $a,b$, $|a-b|\le|a|+|b|$.

\pq[mathematical] Prove that $\big||a|-|b|\big|\le|a-b|$ for all real $a,b$.

\pq[mathematical] Solve $|x-2|>|x|$ and interpret your answer as a statement about distances.

\pq[application] A firm's production line is calibrated to fill bags with $500$ grams of flour. A bag is rejected
if its weight $w$ differs from $500$ grams by more than $2\%$ of $500$ grams. Write the condition for acceptance as an
absolute-value inequality and solve it.

\pq[multi-step] Show that if $|x-3|<1$ then $|x^2-9|<7$. (Hint: write $x^2-9=(x-3)(x+3)$ and bound $|x+3|$.)

\pq[exam-style] (a) State the triangle inequality for three real numbers $x,y,z$ and explain its meaning in terms of
distances on the number line.
(b) Suppose that the real number $c$ satisfies both $|c-1|<\tfrac14$ and $|c|<\tfrac14$. Use the triangle inequality to
derive a contradiction, and explain what this shows about a number that is simultaneously close to $0$ and to $1$.

% --------------------------------------------------------------------
\section{Solutions to practice questions}
% --------------------------------------------------------------------
\pqsol{5.1}
$|-7|=7$; $|3-8|=|-5|=5$; since $\sqrt5>2$, $|2-\sqrt5|=\sqrt5-2$; $|-\tfrac49|=\tfrac49$; $x^2+1\ge1>0$, so
$|x^2+1|=x^2+1$.

\pqsol{5.2}
(a) $-1<x+4<1\Leftrightarrow-5<x<-3$: $x\in(-5,-3)$.
(b) $|3x-6|=3|x-2|\ge9\Leftrightarrow|x-2|\ge3\Leftrightarrow x\le-1$ or $x\ge5$: $x\in(-\infty,-1]\cup[5,\infty)$.
(c) An absolute value is never negative, so $|5-x|\le0\Leftrightarrow|5-x|=0\Leftrightarrow x=5$. Solution $\{5\}$.

\pqsol{5.3}
$a=3$, $b=-5$: $|a+b|=2<8=|a|+|b|$. $a=3$, $b=5$: $|a+b|=8=|a|+|b|$. Equality holds exactly when $a$ and $b$ do not have
opposite signs ($ab\ge0$); when they have opposite signs, part of one cancels part of the other.

\pqsol{5.4}
Apply $|u+v|\le|u|+|v|$ with $u=a$, $v=-b$: $|a-b|\le|a|+|-b|=|a|+|b|$.

\pqsol{5.5}
From $|a|\le|a-b|+|b|$ (\cref{ex:triangle-apply}), $|a|-|b|\le|a-b|$. Exchanging the roles of $a$ and $b$:
$|b|-|a|\le|b-a|=|a-b|$. So both $|a|-|b|$ and its negative are at most $|a-b|$; since $\big||a|-|b|\big|$ equals one of
them, $\big||a|-|b|\big|\le|a-b|$.

\pqsol{5.6}
Both sides are non-negative, so square: $(x-2)^2>x^2\Leftrightarrow x^2-4x+4>x^2\Leftrightarrow-4x>-4\Leftrightarrow x<1$.
Solution $(-\infty,1)$: the points closer to $0$ than to $2$ are those to the left of the midpoint $1$.

\pqsol{5.7}
$2\%$ of $500$ is $10$. Acceptance: $|w-500|\le10\Leftrightarrow490\le w\le510$, i.e.\ $w\in[490,510]$ grams.

\pqsol{5.8}
If $|x-3|<1$ then $2<x<4$, so $5<x+3<7$ and $|x+3|<7$. Then $|x^2-9|=|x-3|\,|x+3|<1\cdot7=7$. (Both factors are
non-negative, so the strict inequalities may be multiplied: if $0\le u<1$ and $0\le v<7$ then $uv\le 1\cdot v<7$ when
$v>0$, and $uv=0<7$ when $v=0$.)

\pqsol{5.9}
(a) For any real $x,y,z$, $|x-z|\le|x-y|+|y-z|$. Since $|u-v|$ is the distance between $u$ and $v$, it says that the
distance from $x$ to $z$ is at most the distance from $x$ to $y$ plus the distance from $y$ to $z$: going via $y$ is
never shorter than going directly.
(b) Take $x=1$, $y=c$, $z=0$:
$1=|1-0|\le|1-c|+|c-0|<\tfrac14+\tfrac14=\tfrac12$, so $1<\tfrac12$, a contradiction. Hence no number can be within
$\tfrac14$ of both $0$ and $1$. This is the key step in showing that $1,0,1,0,\dots$ has no limit (\cref{ex:no-limit}): a
limit would have to be close to both $0$ and $1$ at once.
